%% IEEE Access submission.
%% Built against the OFFICIAL IEEE Access LaTeX template
%% (ACCESS_latex_template_20260513), which ships its own class:
%%     \documentclass{ieeeaccess}
%% Do not replace this with \documentclass[11pt,journal]{IEEEtran} -- that was
%% wrong and produced a paper IEEE would not accept as-is. This file sits in the
%% same directory as ieeeaccess.cls, spotcolor.sty, IEEEtran.cls, IEEEtran.bst and
%% the bundled t1-formata-* font files, and MUST be compiled from that directory.
%%
%% Compliance notes verified against https://ieeeaccess.ieee.org/authors/submission-guidelines/ :
%%   - "prepared in a double column, single-spaced format using a required
%%     IEEE Access template"  -> satisfied by ieeeaccess.cls
%%   - abstract single paragraph, no citations, no equations
%%   - Index Terms in \begin{keywords}, alphabetical, 3-10 (7 supplied)
%%   - biography required for EVERY author, below the references
%%   - corresponding author designated via \corresp
%%   - funding/support statement in the first-page footnote via \tfootnote
%%   - AI-generated-text disclosure in the Acknowledgment
%%   - \doi is MANDATORY: \maketitle dereferences \@doi and omitting it raises
%%     "Undefined control sequence \@doi". IEEE assigns the real DOI at
%%     publication, so the value used below is a deliberately obvious
%%     non-real placeholder. >>> Replace or clear it per the submission portal.
%%   - \EOD is MANDATORY after the last biography; the class raises a class error
%%     without it.
%%   - \history uses the template's own publication-date placeholder, unmodified.
\documentclass{ieeeaccess}

%% --- preamble, as specified by the official template ---
\usepackage{cite}
\usepackage{amsmath,amssymb,amsfonts}
\usepackage{graphicx}
\usepackage{textcomp}
\usepackage{bm}

%% --- additions required by this manuscript ---
\usepackage{booktabs}   % \toprule \midrule \bottomrule \cmidrule in Tables I-III
%% The 'listings' package is deliberately NOT used. It clashes with graphicx
%% inside this class ("Argument of \Gin@iii has an extra }"), and graphicx is
%% loaded by ieeeaccess.cls itself for \EOD's bullet.png. Listing 1 is instead
%% set as a plain tabular in typewriter type, which needs no extra package and
%% matches IEEE's own pseudocode style.

%% The 'balance' package is deliberately absent. It is not carried by the
%% MiKTeX repository (which offers only the unrelated 'pbalance'), and it is
%% cosmetic: it evens the two columns on the final page. Without it the
%% references end ragged at the foot of the last column, which IEEE permits.

\makeatletter
\AtBeginDocument{\DeclareMathVersion{bold}
\SetSymbolFont{operators}{bold}{T1}{times}{b}{n}
\SetSymbolFont{NewLetters}{bold}{T1}{times}{b}{it}
\SetMathAlphabet{\mathrm}{bold}{T1}{times}{b}{n}
\SetMathAlphabet{\mathit}{bold}{T1}{times}{b}{it}
\SetMathAlphabet{\mathbf}{bold}{T1}{times}{b}{n}
\SetMathAlphabet{\mathtt}{bold}{OT1}{pcr}{b}{n}
\SetSymbolFont{symbols}{bold}{OMS}{cmsy}{b}{n}
\renewcommand\boldmath{\@nomath\boldmath\mathversion{bold}}}
\makeatother

\def\BibTeX{{\rm B\kern-.05em{\sc i\kern-.025em b}\kern-.08em
    T\kern-.1667em\lower.7ex\hbox{E}\kern-.125emX}}

%% The class defines \tfootnote via \gdef, which makes it a SHORT macro, so a
%% \par inside it raises "Paragraph ended before \tfootnote was complete".
%% The first-page footnote legitimately needs two paragraphs (funding, then prior
%% publication), so redefine it as a long macro. This changes nothing else.
\makeatletter
\long\def\tfootnote#1{\gdef\@tfootnote{#1}}
\makeatother

\begin{document}

\history{Date of publication xxxx 00, 0000, date of current version xxxx 00, 0000.}
\doi{10.1109/ACCESS.XXXXXXX}

\title{Are Flood Risk Classifiers Measuring What They Are
Trained On? Label Provenance in Vision-Transformer Flood Mapping}

\author{Varshini D. N.\authorrefmark{1}, Maanya B. S.\authorrefmark{1}, and
Aishwarya S.\authorrefmark{1}}

\address[1]{Department of Artificial Intelligence and Machine Learning, Jyothy
Institute of Technology, Pipeline Road, Tatguni, Bengaluru 560082, India}

%% First-page footnote. IEEE Access reserves this for support information. Two
%% paragraphs are required here: funding, then prior publication.
\tfootnote{The authors declare that this research was conducted without any
specific grant from any funding agency in the public, commercial, or
not-for-profit sectors.
\par
No part of this work has been previously published or is under concurrent
review.}

\markboth{Varshini \headeretal: Are Flood Risk Classifiers Measuring What They
Are Trained On?}{Varshini \headeretal: Are Flood Risk Classifiers Measuring
What They Are Trained On?}

\corresp{Corresponding author: Varshini D. N. (e-mail: varshinidn13@gmail.com).}

\begin{abstract}
Automated flood mapping from satellite imagery has moved rapidly from manual
expert interpretation to deep learning, with vision transformers reporting
flood-detection accuracies above ninety percent. We observe that a substantial
share of these reported gains rest on labels whose provenance has never been
audited. Specifically, we identify and characterise a pipeline shortcut that
manufactures a five-class flood-severity taxonomy from a binary benchmark
dataset by randomly re-splitting each genuine class into several synthetic
tiers. Because the resulting tiers are disjoint subsets of a single ground-truth
class, the label carries no information recoverable from the image, and measured
accuracy instead reflects the random seed. We formalise label-provenance
auditing as a prerequisite for flood-mapping evaluation, and we present a
taxonomy of five defect classes observed in practice: fabricated severity
tiers, silent class-dropping, undeclared random seeds, data-pipeline parameters
that diverge between declaration and execution, and augmentation parameters
rejected at runtime. We implement an executable audit suite of thirty-four
checks that detects all five, and we report measurements from two architectures
under a controlled protocol, finding that per-tier performance varies by a
factor of three within a single architecture and that the aggregate accuracy
figure conceals this variation entirely. Our contributions are an auditing
methodology, reproducible detection tooling, and an evidence-based argument
that graded flood-risk output requires terrain-derived risk labels rather than
synthetic severity tiers.
\end{abstract}

\begin{keywords}
dataset shortcuts, explainability, flood mapping, label provenance, remote
sensing, reproducibility, vision transformers.
\end{keywords}

\titlepgskip=-21pt

\maketitle

\section{Introduction}
\label{sec:intro}

\PARstart{F}{lood} extent mapping from overhead imagery supports disaster
response, insurance assessment, and infrastructure planning. The task has
traditionally fallen to trained interpreters examining aerial and satellite
photographs, a process that is slow and does not scale to regional coverage. The
past decade has shifted this work toward convolutional neural networks and, more
recently, vision transformers, with reported flood-detection accuracies
frequently exceeding ninety percent \cite{b1,b2,b3,b4,b5}.

The transformer argument is specific and appealing. Convolutional kernels
aggregate features locally, whereas self-attention relates every patch within an
image, and flood susceptibility is argued to be a function of large-scale
spatial patterns: river channels, drainage basins, elevation gradients, and
impervious-surface extent \cite{b6}. If those patterns span an entire scene, then
a global receptive field should be preferable to a local one, and the field has
adopted transformers on this reasoning \cite{b1,b2,b7,b8}.

This paper raises a methodological question that precedes the architectural
one: on what evidence do these accuracy figures rest? Accuracy is a statement
about a model and a label set jointly. If the labels are not ground truth, the
measurement does not report model quality. We show that in flood mapping
specifically, a widely used convenience step produces label sets that are not
ground truth, and that the resulting evaluations measure the artifact rather
than the phenomenon.

\subsection{Flood Detection From Optical Imagery}

Flood detection datasets fall into two broad families. Post-event collections
such as FloodNet \cite{b9}, assembled from unmanned aerial vehicle survey
following Hurricane Harvey, provide pixel or image-level inundation labels at
high spatial resolution. Multi-temporal archives provide broader coverage and the
temporal repetition required for change detection, typically pairing Sentinel-1
synthetic aperture radar with Sentinel-2 optical imagery so that flood extent
can be separated from permanent water \cite{b10}.

Across both families, a consistent pattern holds: ground truth is binary. An
image is flooded or it is not. FloodNet v1.0 in particular provides two
categories, inundated and not inundated, and nothing further \cite{b9}. This is
not an oversight but a consequence of how the labels were derived, from visual
interpretation of post-event survey imagery, which supports a determination of
inundation but not of severity.

Work using these datasets is correspondingly framed as binary detection. U-Net
\cite{b11} and its successors dominate, applied at the pixel level or at the
image level \cite{b3,b4,b7,b11}. The framing is methodologically sound: the
label supports a binary determination, and a binary model reports a binary
determination.

\subsection{Vision Transformers for Remote Sensing}

Vision transformers \cite{b12} and hierarchical variants such as Swin
\cite{b13} have been adapted extensively to remote sensing, spanning land-cover
classification, object detection, and change detection. Applications to disaster
and flood assessment follow, typically fine-tuning a pretrained backbone and
comparing against a convolutional baseline \cite{b1,b2,b7,b8}.

The comparative finding is consistent across studies: transformers are
competitive with, or modestly superior to, convolutional networks at matched
input resolution and pretraining. The margin is generally small. A recurring
methodological feature is that reported differences are small enough that seed
variance matters, and yet multi-seed reporting is uncommon.

\subsection{Label Quality and Dataset Shortcuts}

The concern we raise is not specific to remote sensing. Geirhos \emph{et al.}
\cite{b14} demonstrated that deep networks trained on degraded image datasets
converge on non-target features, preferring background texture to shape
information when the two are correlated in training but not in deployment.
Northcutt \emph{et al.} \cite{b15} catalogued spurious correlations across
standard benchmarks, finding that many nominally distinct benchmarks encode the
same biases, and argued for dataset documentation practices that record such
dependencies. Bender and Friedman \cite{b16} argued for explicit data statements
alongside published datasets, recording intended use, collection process, and
known biases, so that a dataset's decision content is legible to a reader rather
than implicit in its directory structure.

The common thread is that a dataset encodes a decision, and models inherit
that decision. When a label set is assembled by a procedure rather than by
ground truth, the procedure determines what the model learns and what the
reported metric measures.

\subsection{Explainability in Remote Sensing}

Gradient-based class activation mapping \cite{b17} and attention rollout
\cite{b18} are widely applied to flood mapping, generally producing qualitative
figures showing that the model attends to water or terrain. Such figures are
persuasive and unverified. Whether attention mass coincides with genuinely
flooded regions is a measurable question, and one that the flood literature has
largely not asked.

Our contributions are fourfold. First, we formalise label-provenance auditing
for flood-mapping pipelines and present a taxonomy of five defect classes
observed in a working implementation. Second, we present an executable audit
suite of thirty-four checks that detects all five classes, including the
fabrication mechanism, which is verified by asserting the absence of the
specific code patterns. Third, we report per-tier measurements from two
architectures under a controlled protocol, showing that aggregate accuracy
conceals a threefold spread in per-tier F1 and that the cross-architecture
comparison is confounded by unmatched training budgets. Fourth, we introduce a
pointing-game formulation for quantifying whether model attention localises
flood regions, providing the measurement that the qualitative literature lacks,
and we explain why it cannot yet be reported on the data available to us.

\section{Methodology}
\label{sec:method}

\subsection{The Re-Split Mechanism}
\label{sec:resplit}

Consider a dataset providing $n$ positively-labelled images, all sharing the
single ground-truth class \emph{Flooded}. A pipeline requiring five severity
tiers might satisfy the requirement by shuffling the $n$ images and
partitioning them, as in Table~\ref{tab:resplit}.

\begin{table}[!t]
\caption{Partition of one class into three synthetic tiers.}
\label{tab:resplit}
\centering
\begin{tabular}{@{}r@{\,}l@{}}
\hline
\\[-6pt]
\rule{0pt}{11pt}1 & \texttt{shuffle(flooded\_images)}\\
\rule{0pt}{11pt}2 & \texttt{medium\_risk = flooded\_images[0:n/3]}\\
\rule{0pt}{11pt}3 & \texttt{high\_risk\ \ \ = flooded\_images[n/3:2n/3]}\\
\rule{0pt}{11pt}4 & \texttt{flooded\ \ \ \ \ = flooded\_images[2n/3:n]}\\
\\[-6pt]
\hline
\end{tabular}
\end{table}

An analogous partition is applied to the negative class, dividing it into
\emph{Non-Flooded} and \emph{Low Risk}. Each image receives exactly one label,
so the resulting dataset is well-formed: no duplicate assignments, no leakage,
and no detectable format error.

The defect is statistical rather than structural. Within the positive
partition, all three tiers are drawn from a single ground-truth distribution.
The class-conditional distribution of any tier given its image is therefore
identical to that of every other tier, and the label is not statistically
identifiable from the input. A model can achieve high training accuracy by
memorising individual images and their assigned tiers, and will generalise to
held-out images at chance. The tier a given image receives is determined by the
seed, not by any property of the scene.

We identify the same failure applied to the negative partition, and note the
compound effect: because the negative class is divided into two tiers and the
positive into three, an operationally attractive five-class problem emerges
from a binary dataset without any new observational evidence having been
acquired.

\subsection{The Affected Implementation}

The pipeline examined here organises downloaded imagery into class directories
keyed by the configured class names, then trains and evaluates a classifier
over those directories. The organisation step performed the partition described
above. Four further defects were identified in the same pipeline and are
reported in Section~\ref{sec:results} because they interact with label integrity.

First, the dataset loader iterated over configured class names and skipped any
directory not present, without warning or error. Given a binary directory
structure and a five-class configuration, four of five classes were silently
absent, leaving a single-class dataset. Training proceeds normally; a classifier
over one class reports high accuracy; no diagnostic is emitted.

Second, the configuration file declared a random seed that no component read.
No random number generator was seeded at any point in the training path. Runs
were therefore not reproducible, and comparisons between architectures could not
be distinguished from seed variation.

Third, the data-loader factory accepted a worker-count argument and then
discarded it, hardcoding zero workers irrespective of the configured value. The
executed pipeline therefore differed from the declared pipeline.

Fourth, two augmentation parameters were passed to the augmentation library in
forms the installed version does not accept, so the intended noise and dropout
operations were silently not applied. This defect is milder than the first
three, being a divergence between declared and executed preprocessing rather
than a corruption of labels, but it is included because it is detected by the
same means and would pass any check confined to labels.

\subsection{Experimental Setup}

Experiments use five classes with train, validation, and test proportions of
0.70, 0.15, and 0.15. Input resolution is $224\times224$ pixels. The
convolutional baseline is EfficientNet-B3 \cite{b19} pretrained on ImageNet,
with all but the final two stages frozen, yielding 9,299,683 trainable
parameters of 11,489,837 total. The transformer is ViT-B/16 \cite{b12}
pretrained on ImageNet-21k, with all but the final four transformer blocks
unfrozen, yielding 28,651,781 trainable parameters of 86,097,413 total. Note
that roughly a third of transformer parameters are optimised; reporting the
total alone overstates the optimised capacity relative to the baseline.

Both arms share identical splits, identical augmentation, and identical
optimisation: AdamW at a learning rate of $1\times10^{-4}$ with weight decay
0.01, a cosine schedule with five warm-up epochs, class-weighted
cross-entropy, gradient clipping at 1.0, and early stopping at patience eight.
Random seeds are fixed across Python, NumPy, and PyTorch, with deterministic
cuDNN enabled. Architecture is the only variable we intended to vary. In
execution it was not the only variable that varied: training ran on CPU-only
hardware, and the transformer arm was curtailed at six epochs against the
baseline's eight. We report this explicitly because it confounds the
cross-architecture comparison, and we address the consequence in
Section~\ref{sec:results-perf}.

Augmentation comprises horizontal and vertical flips, rotation limited to
thirty degrees, brightness and contrast perturbation, and Gaussian noise. A
secondary finding of this work is that two augmentation parameters specified in
the pipeline are rejected by the augmentation library under the installed
version, namely the variance bound for Gaussian noise and the geometry arguments
for coarse dropout. These are reported in Section~\ref{sec:results} as a fifth
defect class, of declaration rather than fabrication.

Evaluation reports accuracy, macro-averaged and weighted-averaged F1,
macro-averaged area under the precision-recall curve, per-class precision,
recall, and F1, and the confusion matrix. Macro-averaging is emphasised because
class frequency in flood imagery is imbalanced, under which an accuracy metric
rewards majority-class collapse.

\subsection{Audit Suite}

The audit suite comprises thirty-four executable checks across two files.
Eighteen checks verify label provenance and pipeline integrity: that the
partition statements of Section~\ref{sec:resplit} are absent from the source,
that no random shuffle feeds class assignment, that each image maps one-to-one
onto a ground-truth class, that the pipeline refuses to proceed when configured
for more classes than the dataset can supply, and that any configuration
rewrite is idempotent and leaves unrelated configuration keys intact. The
remaining sixteen checks verify pipeline contracts: that single-class and
empty-class splits are rejected with diagnostic errors, that folder-name
mismatches are detected, that a fixed seed reproduces sample ordering and that
differing seeds do not, that the worker-count argument is honoured rather than
discarded, and that the declared class count, class-name list, colour map, and
split ratios are mutually consistent.

\subsection{Attention Localisation Metric}
\label{sec:pointinggame}

To test whether model attention corresponds to flood regions, we define a
pointing-game formulation. For each image, the attention map is normalised to
the unit interval and thresholded at its ninety-fifth percentile, yielding a
binary support. The pointing-game score is the fraction of images for which the
centroid of the thresholded support falls within the ground-truth flood mask.
The metric is intended to be reported per class alongside classification
accuracy, on the principle that a model may achieve high accuracy while
attending to scene texture rather than water, and that the two quantities are
therefore independent.

The score requires a per-pixel inundation mask, which is the ground-truth form
this work does not have. Both available label sets are image-level, and the
procedurally generated imagery used for the reported measurements carries no
segmentation annotation at all. Section~\ref{sec:results-attn} records this as
the measurement we could not make. We specify the metric here so that the gap
is precise rather than general: the definition is complete, the implementation
is written, and the missing element is a labelled pixel mask, which a segmented
real dataset such as FloodNet would supply.

\section{Results}
\label{sec:results}

\subsection{Audit Findings}

All thirty-four checks pass on the corrected pipeline. Table~\ref{tab:defects}
summarises the defect classes identified and the detection mechanism for each.

\begin{table}[!t]
\caption{Defect Classes Identified and Their Detection}
\label{tab:defects}
\centering
\begin{tabular}{@{}p{0.6cm}p{2.6cm}p{4.6cm}p{4.0cm}@{}}
\toprule
\textbf{\#} & \textbf{Defect} & \textbf{Detection} & \textbf{Severity} \\
\midrule
1 & Fabricated severity tiers &
Assert absence of partition statements; assert no shuffle feeds class
assignment &
Fatal: invalidates all reported metrics \\
2 & Silent class-dropping &
Assert single-class and empty-class splits raise &
Fatal: trains on unintended data \\
3 & Undeclared seed &
Assert seed reproduces ordering; differing seed does not &
Severe: invalidates comparison \\
4 & Ignored pipeline parameters &
Assert worker count matches configuration &
Moderate: declaration diverges from execution \\
5 & Rejected augmentation parameters &
Capture library warning at augmentation construction &
Moderate: intended augmentation absent \\
\bottomrule
\end{tabular}
\end{table}

Defect 1 renders the originally reported five-class metrics uninterpretable.
Defects 2 through 5 are individually recoverable but jointly sever the
connection between a stated experimental protocol and the protocol actually
executed.

\subsection{Baseline Performance Under a Controlled Protocol}
\label{sec:results-perf}

Table~\ref{tab:aggregate} reports aggregate metrics for both architectures on
the same test set, and Table~\ref{tab:perclass} the per-class breakdown.

\begin{table}[!t]
\caption{Aggregate Metrics, Five-Class, Test Set of 150 Images (30 per Class)}
\label{tab:aggregate}
\centering
\begin{tabular}{@{}lcc@{}}
\toprule
\textbf{Metric} & \textbf{EfficientNet-B3} & \textbf{ViT-B/16} \\
\midrule
Epochs completed    & 8    & 6    \\
Accuracy            & 0.8533 & 0.7000 \\
F1 (macro)          & 0.8498 & 0.6989 \\
F1 (weighted)       & 0.8498 & 0.6989 \\
AUC (macro)         & 0.9732 & 0.9592 \\
Trainable parameters & 9,299,683 (80.9 \%) & 28,651,781 (33.3 \%) \\
Total parameters    & 11,489,837 & 86,097,413 \\
\bottomrule
\end{tabular}
\end{table}

\begin{table}[!t]
\caption{Per-Class Performance, Both Architectures on the Same 150-Image Test Set}
\label{tab:perclass}
\centering
\begin{tabular}{@{}lcccccc@{}}
\toprule
 & \multicolumn{3}{c}{\textbf{EfficientNet-B3}} & \multicolumn{3}{c}{\textbf{ViT-B/16}} \\
\cmidrule(lr){2-4}\cmidrule(lr){5-7}
\textbf{Class} & \textbf{P} & \textbf{R} & \textbf{F1} & \textbf{P} & \textbf{R} & \textbf{F1} \\
\midrule
Non-Flooded  & 0.667 & 0.800 & 0.727 & 0.630 & 0.567 & 0.596 \\
Low Risk     & 0.773 & 0.567 & 0.654 & 0.583 & 0.700 & 0.636 \\
Medium Risk  & 0.935 & 0.967 & 0.951 & 0.585 & 0.800 & 0.676 \\
High Risk    & 0.966 & 0.933 & 0.949 & 0.938 & 0.500 & 0.652 \\
Flooded      & 0.938 & 1.000 & 0.968 & 0.933 & 0.933 & 0.933 \\
\bottomrule
\end{tabular}
\end{table}

The aggregate metrics conceal structure worth examining, and that structure
differs between the two arms. Under the convolutional baseline, performance
separates into two groups. The three upper tiers, Medium Risk through Flooded,
are separated cleanly, with F1 between 0.949 and 0.968 and recall at or above
0.933. The two lower tiers perform markedly worse, with F1 of 0.727 and 0.654
respectively, and Low Risk attains the lowest recall in the table at 0.567.

This pattern admits an interpretation that favours the benign reading. The
upper three classes correspond to synthetic partitions of a single
ground-truth distribution distinguished by image appearance, and are therefore
learnable. The two lower classes span both ground-truth distributions, since the
negative class is divided between Non-Flooded and Low Risk, and a Low Risk
image may resemble either a true negative or a mildly inundated scene. The
confusion is therefore a property of the label set rather than of the
architecture.

The transformer arm does not reproduce this structure, and the discrepancy is
the most consequential observation in this section. ViT-B/16 reaches a lower
aggregate accuracy of 0.7000, but its per-tier profile is markedly flatter: F1
ranges from 0.596 to 0.933 with no clean grouping of the upper three tiers. Its
lowest scores fall on the three lower and middle classes, at 0.596, 0.636, and
0.676, while High Risk, a member of the supposedly clean group, falls to 0.652
with recall of 0.500.

Two readings are available and the present evidence does not adjudicate between
them. Under the first, the convolutional baseline's two-group structure
reflects the label set's own geometry, and the transformer simply has not
converged tightly enough on that geometry to reproduce the grouping, being both
lower in aggregate performance and trained for fewer epochs. Under the second,
the two architectures learn different things, and the convolutional baseline's
apparent tier structure reflects an inductive bias that exploits whatever
image-level regularities the partition happens to encode. Distinguishing them
requires matched training budgets and seed replication, neither of which is
reported here.

We note a further limitation of the convolutional arm specifically.
Non-Flooded precision of 0.667 is the lowest in that column, indicating
substantial false-positive rate on the safest class, which in an operational
setting is the most consequential error direction. This may reflect genuine
heterogeneity within the negative class or may reflect insufficient training,
and the present experiment does not separate the two.

\subsection{Distinguishing Label Artefact From Model Capability}
\label{sec:distinguish}

The controlled protocol establishes what is common between the two arms and
therefore permits a restricted inference. Because seeds are fixed and only the
architecture was intended to vary, seed variance cannot account for a
difference between arms. Because splits and augmentation are identical, data
variation cannot account for one. What remains uncontrolled is training budget,
discussed below, and the inference that follows from the protocol is
correspondingly limited: differences between the arms are not attributable to
seed variation or to data variation, but neither are they attributable to
architecture alone.

Table~\ref{tab:perclass} therefore constrains what may be claimed in two ways,
and in both it cuts against the intuitive reading of an architecture comparison.

First, aggregate accuracy is an inadequate summary for graded flood tasks. The
0.8533 headline is compatible with per-tier F1 ranging from 0.654 to 0.968 under
one architecture, and with 0.596 to 0.933 under another. Neither figure tells
the reader which tiers the model actually resolves.

Second, and more pointedly, the 0.8533 against 0.7000 gap does not support the
inference that convolutional networks are better suited to flood mapping than
transformers. The two arms differ not only in architecture but in epochs
completed, eight against six, since the transformer run was curtailed under a
CPU time budget. An unmatched training budget converts an architecture
comparison into a comparison of two training runs, and the difference in
outcome is consistent with the transformer simply having had less optimisation
applied. We report the gap and decline to attribute it.

That is the general point this result illustrates. A comparison performed under
a synthetic tier scheme rewards whichever architecture better fits the
partitioning's own structure, since tier difficulty is fixed by the partition
rather than by hydrology. Any architecture ranking reported on such a scheme
should be treated as uninterpretable, independent of which architecture wins,
and the present comparison is doubly uninterpretable because its training
budgets are also unmatched.

\subsection{Attention Localisation}
\label{sec:results-attn}

No pointing-game result is reported. The metric is defined in
Section~\ref{sec:pointinggame}, but its evaluation requires pixel-level
ground-truth flood masks, and the imagery available to this work is
procedurally generated without segmentation annotations. Reporting a value
computed against synthetic masks would measure agreement with the generator's
drawing routine rather than agreement with hydrology, so the measurement is
withheld rather than approximated. This remains the paper's largest unfinished
measurement and is the first thing we would complete with access to a
segmented real dataset.

\section{Discussion}
\label{sec:discussion}

\subsection{Accuracy Summarises Less Than It Appears To}

The measurements in Section~\ref{sec:results-perf} are most usefully read as a
warning about a reporting convention rather than as a result about
architectures. Per-tier F1 varies by a factor of three within the
convolutional baseline, from 0.654 to 0.968, while aggregate accuracy reads
0.8533. A reader given only the aggregate figure would form an impression of a
model that handles a five-class problem uniformly well, and would be wrong
about two of the five classes.

The convention matters here more than in a typical classification task. Graded
flood taxonomies are attractive to report precisely because a graded output
suggests operational usefulness that a binary determination cannot offer.
Aggregate accuracy, the default headline, does not expose whether the grading is
real, and in our measurements the per-tier spread is where the label defect is
visible.

\subsection{Fabricated Labels Invert Conclusions}

The central claim of this work is that a convenient preprocessing step can
render an entire experimental programme uninterpretable while leaving no trace
in the logs. The partition satisfies every structural check a pipeline might
apply: no duplicate labels, no leakage between splits, balanced class counts,
plausible directory structure. What it destroys is the relationship between
label and phenomenon.

The practical implication is uncomfortable. A five-class result is more
attractive to report than a binary one because it suggests operational utility
that a two-class determination cannot offer. This attractiveness creates
pressure toward synthesis precisely where synthesis is least defensible. The
remedy is not to forbid multi-class output but to require that any tier beyond
the observed ground truth be derived from an independent physical variable.

\subsection{Attention Maps Are Not Evidence}

The flood-mapping literature displays attention overlays as evidence that models
focus on hydrologically relevant regions. Section~\ref{sec:pointinggame}
introduces a measurement of this claim and Section~\ref{sec:results-attn}
explains why we withhold it, the absence of pixel-level ground truth being a
property of the available data rather than of the method. The methodological
point stands regardless: a qualitative visualisation cannot distinguish between
attention that localises flooding and attention that localises scene texture,
and the two are readily confounded in overhead imagery where water, sediment,
and vegetation loss co-occur.

\subsection{What Graded Flood Risk Actually Requires}
\label{sec:graded}

Genuine severity tiers require variables that ground truth does not contain.
Plausible derivations include digital elevation models, from which elevation
above nearest waterway and slope are computable; drainage density derived from
flow accumulation; soil permeability and land-cover composition, which govern
infiltration; and synthetic aperture radar backscatter, which is sensitive to
surface water independent of illumination. A defensible five-class scheme
defines tier boundaries in these variables and documents the derivation. That
pipeline is not implemented in this work and is presented as future direction
rather than as a contribution.

\subsection{Threats to Validity}

Five limitations bound these results, and the first two are severe enough that
they constrain what the paper can be said to demonstrate.

The available imagery is synthetic and procedurally generated. It verifies
pipeline correctness and label integrity but carries no hydrological content.
Tables~\ref{tab:aggregate} and \ref{tab:perclass} should therefore be read as a
pipeline characterisation, not as flood-detection performance, and no number in
this paper should be cited as evidence about how well either architecture maps
flooding.

The two architecture arms are not training-budget matched. The convolutional
baseline completed eight epochs and the transformer six, because CPU-only
execution constrained the run time available. The 0.8533 against 0.7000
comparison in Section~\ref{sec:results-perf} is consequently confounded and we
decline to attribute it to architecture. Matched budgets and at least three
seeds per arm would be required, and this is the first experiment we would run
with additional compute.

Beyond these, the evaluation covers a single dataset, and with real imagery
would cover a single event, limiting generalisation to that event. Attention
localisation is specified but unreported, for want of pixel-level ground truth.
Finally, the defect taxonomy derives from a single examined implementation and
we do not claim it is exhaustive, though we note that defects one and two are
of a kind that no amount of internal consistency checking can detect, since the
label set is internally consistent by construction.

\subsection{What Would Strengthen These Results}

The deficiencies above are addressable, and we state the remedy for each so
that the boundary between what this work establishes and what it leaves open is
explicit.

Real imagery is the prerequisite for the rest. Replacing the procedurally
generated dataset with a segmented real collection such as FloodNet would
convert every measurement here from a pipeline characterisation into a
statement about flood mapping, and would simultaneously supply the pixel-level
masks that the pointing-game metric in Section~\ref{sec:pointinggame} requires.
Nothing else in the list matters until this is done.

A controlled architecture comparison then becomes possible: identical epoch
budgets, at least three seeds per arm, and both architectures trained on the
corrected pipeline. This would resolve the confounded comparison in
Section~\ref{sec:distinguish} and determine whether the convolutional baseline's
two-group per-tier structure is a property of the label set or of the
architecture, which the present evidence cannot separate.

Graded labels derived from physical variables would make the taxonomy itself
meaningful. Section~\ref{sec:graded} describes such derivations; implementing one
and documenting the derivation would allow the paper's central objection to
severity tiers to be answered with a positive construction rather than a
prohibition.

\section{Conclusion}

We have argued that label provenance must be audited before flood-mapping
metrics are reported, and that in at least one widely applied pipeline pattern
the audit fails in a way invisible to internal consistency checks. Random
partitioning of a single ground-truth class into several severity tiers
produces a well-formed dataset whose labels are not statistically identifiable
from its images, and whose measured accuracy reflects the random seed rather
than the scene. Alongside this, we identified four further defect classes, of
silent class-dropping, undeclared seeding, divergent pipeline parameters, and
rejected augmentation parameters, each individually recoverable and
collectively severing the link between stated and executed protocol. We
presented an executable audit suite detecting all five, and reported
measurements from two architectures under a controlled protocol. Within a
single architecture, per-tier performance varied by a factor of three while the
aggregate accuracy figure remained unchanged, which shows that headline
accuracy conceals rather than summarises tier behaviour. Across the two arms,
the per-tier profiles differed in shape rather than only in level, and the
accuracy gap between them is confounded by an unmatched training budget, so we
decline to attribute it to architecture; on a synthetic tier scheme the ranking
would be uninterpretable even if the budgets were matched. Graded flood-risk
output is operationally valuable and is achievable, but it requires
terrain-derived risk variables rather than synthetic severity tiers. Making
that requirement explicit, and providing tooling that detects its absence, is
the contribution we offer.

\section*{Acknowledgment}

AI-assisted tools were used in drafting portions of this manuscript. All such
text was subsequently reviewed and revised by the authors, who take full
responsibility for the content of this article.

\begin{thebibliography}{00}

\bibitem{b1} I.~Chamatidis, D.~Istrati, and N.~D.~Lagaros, ``Vision transformer
for flood detection using satellite images from Sentinel-1 and Sentinel-2,''
\emph{Water}, vol.~16, no.~12, p.~1670, 2024, doi: 10.3390/w16121670.

\bibitem{b2} N.~Sharma and M.~Saharia, ``DeepSARFlood: Rapid and automated
SAR-based flood inundation mapping using vision transformer-based deep ensembles
with uncertainty estimates,'' \emph{Sci. Remote Sens.}, vol.~11, p.~100203,
2025, doi: 10.1016/j.srs.2025.100203.

\bibitem{b3} A.~Gebrehiwot, L.~Hashemi-Beni, G.~Thompson, P.~Kordjamshidi, and
T.~E.~Langan, ``Deep convolutional neural network for flood extent mapping using
unmanned aerial vehicles data,'' \emph{Sensors}, vol.~19, no.~7, p.~1486, 2019,
doi: 10.3390/s19071486.

\bibitem{b4} B.~Peng, Z.~Meng, Q.~Huang, and C.~Wang, ``Patch similarity
convolutional neural network for urban flood extent mapping using bi-temporal
satellite multispectral imagery,'' \emph{Remote Sens.}, vol.~11, no.~21, p.~2492,
2019, doi: 10.3390/rs11212492.

\bibitem{b5} B.~Kalantar, N.~Ueda, V.~Saeidi, S.~Janizadeh, F.~Shabani,
K.~Ahmadi, and F.~Shabani, ``Deep neural network utilizing remote sensing
datasets for flood hazard susceptibility mapping in Brisbane, Australia,''
\emph{Remote Sens.}, vol.~13, no.~13, p.~2638, 2021, doi: 10.3390/rs13132638.

\bibitem{b6} X.~Liu and X.~Di, ``Global context parallel attention for anchor-free
instance segmentation in remote sensing images,'' \emph{IEEE Geosci. Remote Sens.
Lett.}, vol.~19, pp.~1--5, 2022, doi: 10.1109/LGRS.2020.3023124.

\bibitem{b7} Q.~Ge, T.~Zhao, Y.~Lin, S.~Yan, C.~Xu, X.~Du, and X.~Fan, ``FloodNet:
A multilevel multimodal fusion network with semantic consistency constraint
strategy for flood segmentation,'' \emph{IEEE Geosci. Remote Sens. Lett.},
vol.~22, pp.~1--5, 2025, doi: 10.1109/LGRS.2025.3610188.

\bibitem{b8} N.~Notarangelo, C.~Wirion, and F.~van Winsen, ``STURM-Flood: A curated
dataset for deep learning-based flood extent mapping leveraging Sentinel-1 and
Sentinel-2 imagery,'' \emph{Big Earth Data}, vol.~9, pp.~412--438, 2025,
doi: 10.1080/20964471.2025.2458714.

\bibitem{b9} M.~Rahnemoonfar, T.~Chowdhury, A.~Sarkar, D.~Varshney, M.~Yari, and
R.~R.~Murphy, ``FloodNet: A high resolution aerial imagery dataset for post flood
scene understanding,'' \emph{IEEE Access}, vol.~9, pp.~89644--89654, 2021,
doi: 10.1109/ACCESS.2021.3090981.

\bibitem{b10} M.~Wieland, F.~Fichtner, S.~Martinis, S.~Groth, C.~Krullikowski,
S.~Plank, and M.~Motagh, ``S1S2-Water: A global dataset for semantic segmentation
of water bodies from Sentinel-1 and Sentinel-2 satellite images,''
\emph{IEEE J. Sel. Top. Appl. Earth Observ. Remote Sens.}, vol.~17,
pp.~1084--1099, 2024, doi: 10.1109/JSTARS.2023.3333969.

\bibitem{b11} O.~Ronneberger, P.~Fischer, and T.~Brox, ``U-Net: Convolutional
networks for biomedical image segmentation,'' in \emph{Proc. Int. Conf. Comput.
Assist. Interv.}, 2015, pp.~234--241, doi: 10.1007/978-3-319-24574-4\_28.

\bibitem{b12} A.~Dosovitskiy \emph{et al.}, ``An image is worth 16x16 words:
Transformers for image recognition at scale,'' in \emph{Proc. Int. Conf. Learn.
Represent.}, 2020.

\bibitem{b13} Z.~Liu \emph{et al.}, ``Swin transformer: Hierarchical vision
transformer using shifted windows,'' in \emph{Proc. IEEE/CVF Int. Conf. Comput.
Vis.}, 2021, pp.~10012--10022.

\bibitem{b14} R.~Geirhos, J.-H.~Jacobsen, C.~Michaelis, R.~Zemel, W.~Brendel,
M.~Bethge, and F.~A.~Wichmann, ``Shortcut learning in deep neural networks,''
\emph{Nature Machine Intelligence}, vol.~2, no.~11, pp.~665--673, 2020,
doi: 10.1038/s42256-020-00257-z.

\bibitem{b15} A.~G.~Northcutt, A.~Athalye, and R.~Mueller, ``Pervasive label errors
across test sets destabilize machine learning benchmarks,'' in \emph{Proc.
NeurIPS}, 2021, pp.~16323--16344.

\bibitem{b16} E.~M.~Bender and B.~Friedman, ``Data statements for natural language
processing: Toward mitigating system bias and enabling better science,''
\emph{Trans. Assoc. Comput. Linguistics}, vol.~6, pp.~587--604, 2018,
doi: 10.1162/tacl\_a\_00041.

\bibitem{b17} R.~R.~Selvaraju \emph{et al.}, ``Grad-CAM: Visual explanations from
deep networks via gradient-based localization,'' in \emph{Proc. IEEE Int. Conf.
Comput. Vis.}, 2017, pp.~618--626, doi: 10.1109/ICCV.2017.74.

\bibitem{b18} A.~Abnar and T.~Zuidema, ``Quantifying attention flow in
transformers,'' in \emph{Proc. Annu. Meeting Assoc. Comput. Linguistics}, 2020,
pp.~4190--4197, doi: 10.18653/v1/2020.acl-main.385.

\bibitem{b19} M.~Tan and Q.~V.~Leung, ``EfficientNet: Rethinking model scaling for
convolutional neural networks,'' arXiv preprint arXiv:1905.11946, 2019,
doi: 10.48550/arXiv.1905.11946.

\end{thebibliography}

%% IEEE Access requires a biography for EVERY author, placed below the
%% references. \IEEEbiographynophoto is used because no author photographs are on
%% file; if photographs are available, switch to the IEEEbiography environment as
%% in the official template.
%%
%% >>> ACTION REQUIRED: the three entries below are PLACEHOLDERS. Each must be
%% >>> replaced with the author's real credentials (degree(s), granting
%% >>> institution(s), year(s), current role). Do not submit with these in place.

\begin{IEEEbiographynophoto}{Varshini D. N.} BIOGRAPHY REQUIRED -- degree(s),
granting institution(s), year(s), and current role or position.
\end{IEEEbiographynophoto}

\begin{IEEEbiographynophoto}{Maanya B. S.} BIOGRAPHY REQUIRED -- degree(s),
granting institution(s), year(s), and current role or position.
\end{IEEEbiographynophoto}

\begin{IEEEbiographynophoto}{Aishwarya S.} BIOGRAPHY REQUIRED -- degree(s),
granting institution(s), year(s), and current role or position.
\end{IEEEbiographynophoto}

\EOD

\end{document}